\chapter{Blood Alcohol Level (Ethanol)}
\section*{What Is This Test?}
The blood alcohol concentration (BAC) measures the concentration of ethanol (ethyl alcohol) in the blood. Ethanol is a central nervous system depressant that acts primarily by enhancing GABA-A receptor activity (inhibitory) and inhibiting NMDA glutamate receptors (excitatory), producing dose-dependent sedation, disinhibition, impaired judgment, motor incoordination, nystagmus, slurred speech, ataxia, stupor, coma, respiratory depression, and death at very high levels. Ethanol is absorbed rapidly from the stomach and small intestine, distributed throughout total body water, and metabolized primarily by hepatic alcohol dehydrogenase (ADH, 90\%) at a zero-order rate of approximately 10--15 mg/dL/hour (approximately one standard drink per hour) in non-tolerant individuals; chronic alcoholics with CYP2E1 induction metabolize ethanol significantly faster. BAC is used clinically for the evaluation of acute ethanol intoxication, altered mental status of unknown etiology, and suspected concomitant toxic alcohol (methanol, ethylene glycol, isopropanol) ingestion. Legal limits for driving in the United States: 0.08 g/dL (80 mg/dL) for adults $\geq$21 years; 0.02 g/dL (zero tolerance) for drivers <21 years; 0.04 g/dL for commercial vehicle drivers. BAC >300--400 mg/dL (0.30--0.40 g/dL) is potentially lethal from respiratory depression; death most commonly occurs at BAC >400--500 mg/dL (0.40--0.50 g/dL), but chronic alcoholics with tolerance may survive BACs >500--600 mg/dL.
\section*{Alternative Names}
BAC; Blood Alcohol; Ethanol Level; Serum Ethanol; Blood Ethanol Concentration; Alcohol Level
\section*{Principle}
Ethanol is measured by enzymatic assay using alcohol dehydrogenase (ADH) on automated chemistry analyzers. In the presence of ADH and NAD+, ethanol is oxidized to acetaldehyde, and NAD+ is reduced to NADH, which is measured spectrophotometrically at 340 nm. The rate of NADH formation is proportional to ethanol concentration. For forensic/legal purposes, gas chromatography (GC) with flame ionization detection (GC-FID) is used as the confirmatory reference method and is legally defensible.
\section*{History}
The quantitative relationship between alcohol consumption and blood concentration was established by the Swedish physician Erik M.P. Widmark, whose 1932 monograph defined the pharmacokinetic factors (the volume of distribution, r, and the zero-order elimination rate, beta) still used in forensic back-calculation of blood alcohol \cite{jones2010evidence}. In 1927, Emil Bogen measured blood alcohol in traffic accident victims \cite{bogen1927human}, linking alcohol to motor vehicle crashes. Breath alcohol testing evolved in the 1930s--1950s, culminating in Robert Borkenstein's Breathalyzer (1954), which enabled roadside measurement and the landmark Grand Rapids study of drinking drivers (1964). Enzymatic alcohol dehydrogenase assays were adapted to automated analyzers in the 1960s, and gas chromatography became the confirmatory reference method for legal purposes. Legal driving limits were progressively lowered, and the United States adopted a uniform per se limit of 0.08 g/dL (80 mg/dL) by 2004.
\section*{Why This Test Is Ordered}
\begin{itemize}
  \item Evaluation of altered mental status, coma, or suspected acute alcohol intoxication in the emergency department
  \item Medicolegal documentation in driving-under-the-influence cases, workplace accidents, and other forensic settings
  \item Assessment of suspected co-ingestion of toxic alcohols (methanol, ethylene glycol, isopropanol) in patients with an elevated osmolal gap
  \item Monitoring of abstinence in alcohol treatment programs or for patients taking disulfiram
  \item Confirmation of sobriety in specific medicolegal or transplant candidacy assessments
\end{itemize}
\section*{Primary vs Secondary Test}
In the emergency setting, the blood alcohol level is a secondary confirmatory test, ordered after clinical examination and often after an initial roadside breath test. It is the primary quantitative test for forensic and legal purposes, where gas chromatography is the reference method. It is not a screening test for alcohol use disorder.
\section*{How This Test Is Performed}
A blood sample is collected by standard venipuncture, ideally into a sodium fluoride (gray-top) tube, which inhibits in vitro ethanol production by microorganisms; plain serum tubes are acceptable if tested promptly. The venipuncture site must be cleansed with a non-alcoholic antiseptic (e.g., iodine), since alcohol-containing swabs can falsely elevate the result. Ethanol is measured by the enzymatic alcohol dehydrogenase method on automated analyzers for clinical use, with gas chromatography used as the confirmatory reference method for legal specimens. Results are typically available within 1--2 hours for stat clinical samples.
\section*{What Preparation Is Needed}
No fasting is required. For legal specimens, strict chain-of-custody documentation, patient identification, and witnessed collection are essential. The skin should not be cleansed with an alcohol-containing antiseptic. If breath alcohol testing is to be performed, the subject should not have ingested alcohol or alcohol-containing products (including mouthwash) within the preceding 15 minutes, but this does not apply to the blood test.
\section*{Contraindications and Precautions}
There are no absolute contraindications to venipuncture. Precautions: the venipuncture site must be prepared with a non-alcoholic antiseptic to avoid sample contamination; a fluoride-containing gray-top tube should be used when testing may be delayed, as ethanol can be produced in vitro by contaminating organisms in plain tubes; a single result reflects one point in time, and back-calculation of a prior blood alcohol is unreliable when the exact time of ingestion is unknown; and endogenous ethanol production (auto-brewery syndrome) is an extremely rare but recognized cause of a detectable level without ingestion.
\section*{Risks}
Minimal risk. Slight pain or bruising at the venipuncture site. Rarely: hematoma, infection, or vasovagal syncope.
\section*{What Happens After the Test}
Results return within 1--2 hours and are interpreted against legal limits (0.08 g/dL for most adult drivers in the United States) and the clinical examination. Management of acute intoxication is supportive, with airway protection, monitoring of blood glucose, and treatment of respiratory depression in severe cases. If altered mental status does not resolve as the blood alcohol level falls, co-ingestion of toxic alcohols (methanol or ethylene glycol) and other causes must be investigated with serum osmolality and the osmolal gap.
\section*{Understanding Results}
\subsection*{Not Detected (<10 mg/dL)}
No significant recent ethanol ingestion. Does not exclude chronic alcoholism or alcohol use disorder.
\subsection*{Low (10--50 mg/dL, 0.01--0.05 g/dL)}
Mild impairment in non-tolerant individuals. Euphoria, mild disinhibition.
\subsection*{Moderate (50--150 mg/dL, 0.05--0.15 g/dL)}
Ataxia, slurred speech, nystagmus, impaired judgment, prolonged reaction time. Most U.S. states' legal definition of driving under the influence (DUI) is $\geq$80 mg/dL (0.08 g/dL).
\subsection*{High (150--300 mg/dL, 0.15--0.30 g/dL)}
Vomiting, severe ataxia, confusion, stupor. May progress to coma.
\subsection*{Potentially Lethal (300--400 mg/dL+, 0.30--0.40+ g/dL)}
Coma, respiratory depression, hypothermia, hypotension, death. Can be higher in chronic alcoholics with tolerance.
\section*{Normal Results}
Negative (<10 mg/dL, or 0.00 g/dL). No detectable ethanol. Legal DUI: $\geq$80 mg/dL (0.08 g/dL). Not legally intoxicated: <80 mg/dL (but still potentially impaired; DUI can be charged at lower levels + behavioral evidence). Endogenous ethanol production (fermentation of carbohydrates by gut bacteria, ``auto-brewery syndrome'') is an extremely rare condition producing detectable ethanol without ingestion.
\section*{Follow-up Tests}
Serum osmolality and osmolal gap (for co-ingestion of toxic alcohols --- methanol, ethylene glycol, isopropanol), basic metabolic panel, serum glucose, and if altered mental status persists after ethanol has been metabolized: CT head, toxicology screens.
\section*{References}
\bibliographystyle{unsrtnat}
\bibliography{references}
\clearpage
\section*{Regulatory Status and Authorization Date}
Regulatory status applies to the specific assay, instrument, manufacturer, and intended use rather than to the test name alone. No single approval date can be assigned to this test category. For a commercial assay, verify the current product in the relevant regulator's database: the U.S. Food and Drug Administration (FDA) clearance, approval, or authorization; India's Central Drugs Standard Control Organisation (CDSCO) licence or permission; the European Union In Vitro Diagnostic Regulation (EU IVDR) CE certificate issued through the applicable conformity-assessment route; or the United Kingdom Medicines and Healthcare products Regulatory Agency (MHRA) registration and UKCA/CE status. Laboratory-developed tests may not have an individual FDA, CDSCO, EU IVDR, or MHRA approval date.
\clearpage
